\section{Architektura modelu}
\label{sec:arch}

\subsection{Reprezentacja tekstu: TF-IDF}

Tekst jest zamieniany na wektor cech metodą \textbf{TF-IDF} (\emph{Term
Frequency--Inverse Document Frequency}). Jest to nieparametryczna inżynieria cech ---
\emph{nie} model pretrenowany --- więc model neuronowy w całości uczy się od zera.
Wektoryzator (\texttt{make\_tfidf\_vectorizer}, \texttt{src/model/baseline.py}) jest
\textbf{współdzielony} z baselinem (regresją logistyczną), co zapewnia uczciwe
porównanie obu klasyfikatorów. Jest dopasowywany \textbf{wyłącznie na zbiorze
treningowym}, a tym samym słownikiem transformowane są val i test (brak wycieku
ze zbiorów ewaluacyjnych). Parametry zebrano w Tabeli~\ref{tab:tfidf}.

\begin{table}[H]
  \centering
  \caption{Konfiguracja wektoryzatora TF-IDF.}
  \label{tab:tfidf}
  \begin{tabular}{ll}
    \toprule
    \textbf{Parametr} & \textbf{Wartość} \\
    \midrule
    \texttt{max\_features} & 50\,000 (rozmiar słownika) \\
    \texttt{ngram\_range}  & (1, 2) --- unigramy i bigramy \\
    \texttt{min\_df}       & 5 (minimalna częstość dokumentowa) \\
    \texttt{sublinear\_tf} & \texttt{True} ($1 + \log(\mathrm{tf})$) \\
    \texttt{strip\_accents}& \texttt{"unicode"} \\
    \bottomrule
  \end{tabular}
\end{table}

\subsection{Sieć gęsta (MLP)}

Klasyfikatorem jest wielowarstwowa sieć gęsta (\texttt{MLPClassifier},
\texttt{src/model/mlp.py}) o architekturze:
\[
  \underbrace{\text{Linear}(50\,000, 256)}_{\text{warstwa ukryta}}
  \;\rightarrow\; \text{ReLU}
  \;\rightarrow\; \text{Dropout}(0{,}3)
  \;\rightarrow\; \underbrace{\text{Linear}(256, 1)}_{\text{wyjście (logit)}}.
\]
Wyjściem jest pojedynczy \emph{logit} (bez sigmoidy), co pozwala użyć numerycznie
stabilnej funkcji straty \texttt{BCEWithLogitsLoss}. Sieć liczy
\textbf{12\,800\,513} parametrów --- zdecydowana większość przypada na pierwszą
warstwę ($50\,000 \times 256$), bo wymiar wejścia równy jest rozmiarowi słownika.

\subsection{Proces uczenia}

Funkcja \texttt{train\_mlp} realizuje pełną pętlę treningową z \textbf{early
stopping} monitorującym F1 na zbiorze walidacyjnym i przywracającym wagi z najlepszej
epoki (Tabela~\ref{tab:train}). W praktyce trening zatrzymał się po \textbf{11
epokach} (z dozwolonych 20).

\begin{table}[H]
  \centering
  \caption{Hiperparametry uczenia.}
  \label{tab:train}
  \begin{tabular}{ll}
    \toprule
    \textbf{Parametr} & \textbf{Wartość} \\
    \midrule
    optymalizator        & Adam \\
    learning rate        & $1\times10^{-3}$ \\
    funkcja straty       & \texttt{BCEWithLogitsLoss} \\
    rozmiar batcha        & 256 \\
    maks. epok           & 20 \\
    early stopping       & na F1 walidacji, \texttt{patience = 3} \\
    dropout              & 0{,}3 \\
    weight decay (L2)    & 0 (domyślnie; badany w sekcji~\ref{sec:eval}) \\
    próg decyzyjny       & 0{,}5 \\
    ziarno losowe        & \texttt{RANDOM\_STATE} (powtarzalność) \\
    \bottomrule
  \end{tabular}
\end{table}

\paragraph{Uzasadnienie wyborów.}
Funkcja \texttt{BCEWithLogitsLoss} łączy sigmoidę z entropią krzyżową w sposób
numerycznie stabilny i jest naturalna dla klasyfikacji binarnej. Adam zapewnia szybką
zbieżność bez ręcznego strojenia harmonogramu uczenia. Early stopping na F1 (a nie na
stracie) jest spójny z metryką raportowaną i chroni przed przeuczeniem przy
praktycznie zerowym koszcie. Dropout 0{,}3 to standardowa, umiarkowana regularyzacja
warstwy ukrytej.
